\documentclass{article}
\usepackage[T1]{fontenc}
\usepackage{lmodern}
\usepackage{graphicx}
\usepackage{amsmath,amssymb}
\usepackage[a4paper,margin=2cm]{geometry}
\usepackage[absolute,overlay]{textpos}
\setlength{\TPHorizModule}{1mm}
\setlength{\TPVertModule}{1mm}
\usepackage[indonesian]{babel}
\usepackage{hyperref}
\setlength{\emergencystretch}{2em}
% unit: Halaman 2

\begin{document}

% === Halaman 2 (llm) ===

\section*{Pendahuluan}

\begin{itemize}
    \item Kriptografi modern adalah era kriptografi setelah penemuan komputer digital (awal tahun 1970-an).
    \item Perkembangan teknologi komputer digital membuat hampir semua bidang ilmu pengetahuan berkembang pesat, termasuk kriptografi.
    \item Komputer digital merepresentasikan data dan informasi dalam biner.
    \item Algoritma kriptografi modern beroperasi dalam mode bit (bandingkan dengan algoritma kriptografi klasik beroperasi dalam mode karakter)
    \begin{itemize}
        \item kunci, plainteks, cipherteks, diproses dalam rangkaian bit
        \item operasi \textbf{xor} paling banyak digunakan di dalam algoritmanya
    \end{itemize}
\end{itemize}

\end{document}
